\documentclass[11pt]{article}
\usepackage[T1]{fontenc}
\usepackage{amsmath,amssymb,amsthm,mathtools}
\usepackage{graphicx}
\usepackage[margin=0.9in]{geometry}
\usepackage{microtype}
\usepackage[hidelinks]{hyperref}

\newtheorem{theorem}{Theorem}
\newtheorem{remark}[theorem]{Remark}
\newcommand{\R}{\mathbb R}

\title{A smooth-measure counterexample to the direct generalization of Theorem 3.4\\
\large A scoped negative answer to Question 6.1 of arXiv:2510.00645}
\author{Autonomous proof packet --- agent\_lane\_12}
\date{11 August 2026}

\begin{document}
\maketitle

\begin{abstract}
Malliaris and Mar\'in Sola ask whether their weighted lower bound and a
corresponding upper bound extend to general measures, concavity transforms,
and functions $N$, explicitly highlighting discrete measures.  We show that
the direct lower-bound extension is false, even while retaining
$\psi=\log$ and the affine function $N(t)=t$.  An exact three-atom example
already violates the formula after the most natural fractional-atom repair.
Replacing the atoms by three narrow translates of one smooth bump preserves
the strict violation.  Thus there is a counterexample whose measure has a
compactly supported $C^\infty$ density, so neither atom conventions nor a
discontinuous distribution function cause the failure.  Alternative
measure-dependent formulations and the upper-bound part remain open.
\end{abstract}

\section{Source question and the direct extension}

The source is A.~Malliaris and F.~Mar\'in Sola, \emph{Quantitative
improvements of functional inequalities under concavity properties},
arXiv:2510.00645~\cite{source}.  Question 6.1 on PDF page 17 asks whether
Theorem 3.4 and the corresponding upper bound hold for general $\mu,\psi,N$;
the following paragraph specifically names discrete $\mu$ as an untreated
case.

\begin{center}
\includegraphics[width=0.98\textwidth]{figures/open_problem_crop.png}
\end{center}

For the measures treated in Theorem 3.4, let
\[
 \Phi(t)=\mu([0,t]),\qquad
 u=\int_0^h f\,d\mu,\qquad V=\int_0^\infty f\,d\mu.
\]
The theorem states, for increasing convex $N$, decreasing log-concave $f$,
and $u\leq V/2$, that
\begin{equation}\label{eq:extension}
 \int_0^\infty N(t)f(t)\,d\mu(t)
 \geq \frac{u}{\Phi(h)}
 \int_0^{\Phi^{-1}((V/u)\Phi(h))}N(t)\,d\mu(t).
\end{equation}
The direct interpretation of the lower-bound part of Question 6.1 is to
retain \eqref{eq:extension} for an arbitrary positive Borel measure.  The
counterexample below has a continuous distribution function and a unique
relevant quantile, so the inverse in \eqref{eq:extension} is unambiguous.

\section{Proof intuition}

The proposed extremizer effectively fills the first $(V/u)\Phi(h)$ units of
$\mu$-mass at the average height $u/\Phi(h)$.  A general measure can place a
large block of mass just to the right of $h$.  A decreasing exponential $f$
discounts that block, while the increasing $N$ magnifies its location.  With
three carefully weighted blocks, the total-mass constraint is exact but the
putative filled-prefix moment is too large.

The atomic calculation has a gap of $1/4$.  Smoothing each atom by the same
narrow even bump preserves the exponential mass ratios exactly.  The only
errors are first moments inside intervals of radius $\varepsilon$, and their
total size is less than the atomic gap when $\varepsilon<1/18$.

\section{Counterexample with a smooth density}

\begin{theorem}[Failure for smooth positive Borel measures]\label{thm:main}
There exist a finite positive Borel measure $\mu$ on $[0,\infty)$ with a
compactly supported $C^\infty$ density, a decreasing integrable log-concave
function $f$, an increasing convex function $N$, and $h>0$ such that
$u=V/2$ but the inequality \eqref{eq:extension} is strictly reversed.
Consequently the direct lower-bound extension of Theorem 3.4 to general
measures is false, already for $\psi=\log$ and affine $N$.
\end{theorem}

\begin{proof}
Fix $0<\varepsilon<1/18$, and let $\rho\in C_c^\infty(\R)$ be an even
nonnegative probability density supported in $(-\varepsilon,\varepsilon)$.
On $[0,\infty)$ define
\[
 d\mu(t)=m(t)\,dt,
 \qquad
 m(t)=\rho(t-1)+\rho(t-2)+6\rho(t-3).
\]
Take
\[
 f(t)=2^{1-t},\qquad N(t)=t,\qquad h=\frac52.
\]
The supports of the three bumps are disjoint and lie in $(0,\infty)$.
Clearly $f$ is decreasing, Lebesgue integrable, and log-concave (indeed,
$\log f$ is affine), while $N$ is nonnegative, increasing, and convex.

Put
\[
 A=\int_{-\varepsilon}^{\varepsilon}2^{-s}\rho(s)\,ds>0,
 \qquad
 B=\int_{-\varepsilon}^{\varepsilon}s,2^{-s}\rho(s)\,ds.
\]
Since $|s|<\varepsilon$ on the support, $|B|\leq\varepsilon A$.
Translation of the same bump and the identity
$f(j+s)=2^{1-j}2^{-s}$ give
\[
 u=\int_0^h f\,d\mu=A+\frac A2=\frac{3A}{2},
 \qquad
 V=A+\frac A2+6\frac A4=3A=2u.
\]
Also $\Phi(h)=2$, since precisely the first two unit-mass bumps lie to the
left of $h$.  Hence the mass level in \eqref{eq:extension} is
\[
 \frac Vu\Phi(h)=4.
\]

Let $d=\Phi^{-1}(4)$.  The point $d$ lies inside the third bump, and if
$q=d-3$, then
\[
 \int_{-\varepsilon}^{q}\rho(s)\,ds=\frac13.
\]
Writing $C=\int_{-\varepsilon}^{q}s\rho(s)\,ds$, we have
$C\geq-\varepsilon/3$.  Evenness of $\rho$ makes its full first moment zero,
so the prefix $N$-moment is
\begin{align*}
 Q:=\int_0^d t\,d\mu(t)
 &=1+2+6\int_{-\varepsilon}^{q}(3+s)\rho(s)\,ds\\
 &=9+6C\geq9-2\varepsilon.                         \tag{1}
\end{align*}
Thus the proposed right-hand side satisfies
\begin{equation}\label{eq:rhs}
 \frac{u}{\Phi(h)}Q
 =\frac{3A}{4}Q
 \geq \frac{27A}{4}-\frac{3A}{2}\varepsilon.
\end{equation}

The actual left-hand side can be evaluated on the three translated bumps:
\begin{align*}
 \int_0^\infty t f(t)\,d\mu(t)
 &=(A+B)+\frac12(2A+B)+\frac{6}{4}(3A+B)\\
 &=\frac{13A}{2}+3B
 \leq \frac{13A}{2}+3\varepsilon A.                \tag{2}
\end{align*}
Subtracting the upper bound in (2) from the lower bound in
\eqref{eq:rhs} gives
\[
 A\left(\frac14-\frac92\varepsilon\right)>0.
\]
Therefore the right-hand side of \eqref{eq:extension} is strictly larger
than its left-hand side, proving the claimed failure.
\end{proof}

\section{Exact atomic core and robustness}

The numerical choices become transparent before smoothing.  Let
\[
 \mu_0=\delta_1+\delta_2+6\delta_3,
 \qquad f(t)=2^{1-t},\quad N(t)=t,quad 2<h<3.
\]
Then $\Phi(h)=2$, $u=3/2$, $V=3$, and
\[
 \int Nf\,d\mu_0=1+2\left(\frac12\right)
       +6\cdot3\left(\frac14\right)=\frac{13}{2}.
\]
The standard generalized inverse gives $\Phi^{-1}(4)=3$ and includes the
whole atom at $3$, yielding the much larger proposed right side $63/4$.
More importantly, even if that atom is split fractionally, the first four
units of measure have $N$-moment
\[
 1+2+2\cdot3=9,
\]
so the repaired right side is $27/4>13/2$ by exactly $1/4$.

There is also a one-parameter family.  For $0<r<1$, replace $f$ by
$r^{t-1}$ and the last weight $6$ by $(1+r)/r^2$.  Again $V=2u$, and the
fractional-prefix right side minus the actual moment is exactly
$(1-r)/2>0$.  The selected value $r=1/2$ merely keeps every quantity simple.

\section{Verification, scope, and novelty guardrail}

The script \texttt{code/verify\_counterexample.py} uses rational arithmetic
to check the atomic example, the fractional-atom gap, the entire rational
parameter family sampled at $r=k/100$, and the rigorous smooth-error margin
at $\varepsilon=1/100$.  It reports that all checks pass.  The smooth proof
itself is symbolic and does not depend on computation.

This disproves only the natural same-form lower bound.  It neither rules out
a regularity- or measure-dependent correction nor settles the upper-bound
questions.  The smooth density shows that excluding atoms is insufficient.

On 11 August 2026, bounded novelty checks covered the four run indexes and
arXiv/web searches by exact title, id, Question 6.1, authors, and
``general/discrete measure.''  They found the source but no later answer or
matching counterexample; this is not a priority claim.

\paragraph{Human-review recommendation.}
Accept as a scoped counterexample to the direct lower-bound half of Question
6.1.  Check that retaining \eqref{eq:extension} is the intended
interpretation; all hypotheses and the strict gap are then explicit.

\begin{thebibliography}{9}
\bibitem{source}
A.~Malliaris and F.~Mar\'in Sola,
\emph{Quantitative improvements of functional inequalities under concavity
properties}, arXiv:2510.00645 (2025), Question 6.1 and Theorem 3.4.
\end{thebibliography}

\end{document}
